\section{Discussion}
\label{sec:discussion}

\subsection{Internal versus external behaviour}
With patient- and duplicate-grouped splits and no proxy labels, a standard DenseNet-121 reached an internal test accuracy of \RTestAcc\% and a macro ROC-AUC of \RTestMacroAUC. On radiographs from a hospital it had never seen, the same model flagged \RExtSensK/\RExtSensN{} TB scans as abnormal but called only \RExtSpecK/\RExtSpecN{} normal scans Normal, and it named \RExtTBK/\RExtTBN{} TB scans Tuberculosis, although its internal Tuberculosis recall was \RRecallTB\%. Its external predictions were dominated by Pneumonia and Nodule/Mass, classes the cohort does not contain. The internal results therefore overstate the behaviour of this model on an unseen hospital by a wide margin, and they do so even though the usual leakage paths were closed: leakage control is necessary for a credible internal estimate, but it does not make that estimate a measure of generalisation.

Four observations point to source confounding as one reason, and we stress that each is indirect. First, the benchmark is confounded by design: classes come from different subsets of sources, and the source of an image can be predicted from frozen ImageNet features with \RProbeImgBalAcc\% balanced accuracy (chance \RProbeImgChance\%), and still within a single class, for example \RProbeImgWithinNormal\% among Normal images (chance \RProbeImgWithinNormalChance\%). Source information is thus linearly available in the features on which the classifier's head operates. Second, internal accuracy varies widely between sources (Table~\ref{tab:persource}), and Normal images were recognised much less often when they came from NIH than from other sources. Third, run~1, whose only COVID-19 source was the COVID-19 image data collection, called about half of the Montgomery radiographs COVID-19 (\RRunOneExtCovidCalls{} of \RNExternal);
 when build~2 added a large second COVID-19 source, run~2 almost stopped doing so, which is what one would expect if run~1 had linked the appearance of that single source with the class. Fourth, the pre-specified Grad-CAM examples for COVID-19 and Normal highlight the neck, shoulders and image borders rather than the lungs. None of these shows that the classifier relies on source cues; together they make that explanation plausible, and a causal test (leave-one-source-out training, masking of regions outside the lungs) remains to be done.

Adding sources did not close the external gap. Run~2's errors on Montgomery moved from COVID-19 to Pneumonia and Nodule/Mass instead of becoming correct, its external specificity was lower than run~1's, and Tuberculosis was still not recognised although it now comes from two sources (Shenzhen and TBX11K). Adding a second source per class changed where the external errors went rather than removing the external gap. On the common selection set run~2 was clearly better, but most of that advantage lies on the sources only run~2 trained on: on the four shared sources the two runs had nearly the same accuracy (Section~\ref{sec:res-runs}). These are single-seed observations; Section~\ref{sec:res-seeds} reports how much seeds alone change the results.

Finally, the external ROC-AUC of the abnormality score (\RExtAUC) is well above chance while specificity is low, just as the internal macro ROC-AUC is much higher than internal accuracy. Ranking metrics reward scores that order positives above negatives, whatever the absolute level; on a new hospital, a shift of the scores towards ``abnormal'' can leave the ranking partly intact and still move most normal images across the arg-max boundary. A model that is to be deployed at a new site would at least need its operating point re-established there, and a high ranking metric on internal data is not evidence that this is unnecessary. We do not compare these numbers with those of other studies, which used different data and splits.

Independently of the final numbers, the design of the evaluation determines what can be concluded. Internal test images share sources, acquisition devices, post-processing and, for some sources, populations with the training images. For a class that is supplied by one collection, the internal test set cannot distinguish a model that recognises the disease from one that recognises the collection~\cite{zech2018variable,degrave2021ai,maguolo2020critic}. The external cohort breaks this link for two classes only (Normal and Tuberculosis), and it does so in a deliberately weak way: it asks whether normal radiographs from an unseen hospital are called Normal and whether TB radiographs are called abnormal and, more strictly, Tuberculosis. Reporting the full table of predicted classes on this cohort, rather than a single sensitivity, is what makes a shortcut visible: an abnormal prediction counts towards sensitivity even if the predicted disease is one the cohort does not contain.

\subsection{What the audit of the earlier pipeline showed}
\label{sec:disc-audit}
The V5 pipeline that preceded this work is the clearest illustration of why these controls matter. V5 pooled eight public sources, several of them obtained from third-party redistributions, into 10,547 images, and its partitions were built incrementally across dataset versions. Its canonical model recorded an internal test accuracy of \RVFiveAcc\% on $N=\RVFiveN$ images. A retrospective audit of the manifests, ingestion code and recorded results found the following.
\begin{itemize}
\item \textbf{Patient leakage.} Patient tokens recoverable from Kermany file names showed that 154 of the 1,570 internal test images (9.8\%) belonged to patients who also had images in the training split. A further 8,478 images had no real patient identity, so their separation could not be checked at all.
\item \textbf{Proxy and latent labels.} 1,395 Pneumonia labels (about half of the class) were TBX11K images labelled ``sick but non-TB'', whose cause is not specified; 124 Tuberculosis labels were latent TB; and 1,440 of 1,467 VinDr-CXR images used as single-label examples carried two or more findings.
\item \textbf{Degenerate external behaviour.} On the Montgomery cohort, every recorded V5 variant (five training paradigms, including two domain-adaptation methods) flagged essentially every TB radiograph as abnormal but called none of the normal radiographs Normal (specificity \RVFiveExtSpecK/\RVFiveExtSpecN). A sensitivity quoted alone would have suggested a perfect screening tool; with specificity it describes a model that never outputs Normal on that hospital.
\item \textbf{Possible contamination of the external cohort.} One V5 TB source is described as containing Montgomery images. A content check (MD5 and perceptual hash) found no match, but it also found no match for Shenzhen images that the archive is described as containing, so the check cannot rule the overlap out. V6 does not use that archive.
\item \textbf{Reporting errors.} A comparison table in earlier drafts differed from the result files in 13 of 35 cells, all in the direction of making comparator runs look worse, and three figures captioned as the canonical model showed a baseline run. The canonical checkpoint itself was lost, so ROC-AUC and calibration could not be re-verified.
\end{itemize}
None of these defects is unusual in isolation; together they produced an internal accuracy that said little about the model. V6 was designed so that each defect is either removed (proxy labels, aggregated archives, per-image splitting of patients with known identity) or explicitly measured and reported (unknown identity, near-duplicates, external specificity, per-source performance). The V5 numbers are not comparable with V6: the data, the splits and the label definitions differ.

\subsection{Implications for multi-source benchmark design}
From the audit and the rebuild we draw the following recommendations, which do not depend on the final V6 numbers.
\begin{enumerate}
\item \textbf{Acquire from official sources and log provenance.} Third-party redistributions rename files, drop patient identifiers and silently merge collections; the overlap between a redistributed TB archive and our external cohort could not be resolved for this reason. A per-file record of URL, size and SHA-256 makes a dataset rebuildable without redistributing images.
\item \textbf{Fix labelling rules before training, and label from the source's own labels only.} Rules that map ``abnormal, not TB'' to ``pneumonia'' or latent TB to TB change the meaning of a class; they should be excluded, and the exclusion counted.
\item \textbf{Split by patient and by content.} Grouping by the genuine patient identifier where one exists, by a file-name-derived identifier where possible, and merging exact and near-duplicate images with a two-stage test (perceptual hash and thumbnail correlation) closes the most common leakage paths. Where identity is unknown, say so and report how many images are affected; zero groups spanning splits is not proof of patient separation for those images.
\item \textbf{Report per source.} Overall accuracy is a weighted mixture of per-source accuracies and changes with the source mix; a class supplied by one source should be flagged as confounded.
\item \textbf{Evaluate on an unseen hospital and always report specificity with sensitivity}, together with the exact-class recall and the full prediction breakdown.
\item \textbf{Pre-register selection rules.} Choosing between runs or builds on a set that neither has trained on, with the rule written before results exist, prevents test-set selection.
\end{enumerate}

\subsection{From benchmark to deployment: the input gate}
\label{sec:disc-gate}
The classifier is served by a web application that refuses images judged not to be chest radiographs. The intended gate was a trained model; because its weights were not available, the application fell back to a heuristic score built from image size, aspect ratio, contrast, colour, brightness and an anatomical intensity profile, with the decision threshold (0.83) inherited from the trained gate. We measured this gate on the development data. At that threshold it rejected 1,096 of 1,724 genuine radiographs of the build-1 test split (63.6\%; from about 34\% of Shenzhen images to 83\% of COVID-19 images from the COVID-19 image data collection) and 83 of the 138 Montgomery radiographs (60.1\%).
% Source: artifacts/app_checks/cxr_gate_rejection_v6_build_001.json (application check, not a model result).
In practice, the application refused most of the radiographs it was built for, while the model metrics, which are computed without the gate, gave no hint of this. We recalibrated the heuristic threshold on the build-2 validation radiographs only, as the larger of the threshold that passes 98\% of them and the highest score of any synthetic non-radiograph that passed the first-stage checks plus 0.01, giving 0.703. At this threshold 91.0\% of the validation radiographs pass (38.7\% at 0.83), and all seven synthetic non-radiograph test images (standing in for a colour and a grey-level photograph of a car, a dog, a selfie, a landscape, a screenshot and a computed tomography slice) are rejected.
% Source: artifacts/app_checks/cxr_gate_calibration_v6_build_002.json.
Re-measured with the recalibrated threshold, the gate rejected 218 of the 2,350 radiographs of the build-2 test split (9.3\%) and 18 of the 138 Montgomery radiographs (13.0\%), against 63.6\% and 60.1\% before.
% Source: artifacts/app_checks/cxr_gate_rejection_v6_build_002_heuristic0703.json (application check, not a model result).
% Note: the two "before" figures were measured on the build-1 test split, the "after" figure on the build-2 test split.
The before and after test figures refer to different test splits (build~1 and build~2), and the threshold was set on validation radiographs only, so the test measurement is a check, not a tuning target. The negatives used for calibration are synthetic and few, so the gate's ability to reject real non-radiographs is not established; about one genuine radiograph in ten is still refused. The general point is that a component outside the classifier can dominate the behaviour of a deployed system, and it must be evaluated on the same data as the model.

\subsection{Limitations}
\label{sec:limitations}
\begin{itemize}
\item \textbf{Few runs, CPU, time-boxed.} Runs 1 and 2 are single training runs (seed \RSeed) on a laptop CPU, bounded by an epoch limit or a wall-clock budget. Seed variability is estimated from three runs of the selected configuration only (Section~\ref{sec:res-seeds}), which gives a rough spread, not a reliable variance; the run comparison has one seed per run. The bootstrap and Wilson intervals describe test-set sampling only.
\item \textbf{Limited ablations.} The preprocessing chain was taken from the earlier V5 work, where the low-pass filter was chosen on V5 validation data; on V6 only the full chain and no filtering are compared, in one run each (Section~\ref{sec:res-ablation}). Single filters, the filter order relative to resizing, source composition and alternative backbones were not studied.
\item \textbf{Shortcut evidence is indirect.} The source-identifiability probe shows that source information is available in the features, not that the classifier uses it; together with the per-source, external and Grad-CAM results it supports a source shortcut as a plausible explanation, not a demonstrated one. No leave-one-source-out training or masking experiment was carried out.
\item \textbf{One small, binary external cohort.} Montgomery has 138 images and two classes. It cannot test four of the six classes, and its confidence intervals are wide.
\item \textbf{External cohort not strictly untouched.} Run~2 was trained after run~1's external results had been seen (Section~\ref{sec:design}); a fresh external cohort is needed for a strictly prospective evaluation.
\item \textbf{Unknown patient identity} for TBX11K and the COVID-19 Radiography Database (build~2): these images are grouped per image plus merged near-duplicates, so the same patient could appear in two splits.
\item \textbf{Label quality and the single-label assumption.} NIH labels are mined from reports, not read from images. Images with several findings were excluded, so the test set contains only clear single-label cases, unlike routine practice.
\item \textbf{Population and source confounding.} Kermany images are paediatric, all other sources adult; no demographic data were available, so subgroup performance~\cite{seyyedkalantari2021underdiagnosis} could not be assessed. COVID-19 and Tuberculosis come from one source each in build~1 and two in build~2. Pleural Effusion and Nodule/Mass come from NIH only because VinDr-CXR and JSRT were not accessible.
\item \textbf{Resolution-dependent preprocessing.} Filters run at native resolution before resizing, so images of different resolution receive different effective filtering. Training and serving use the same code, so they are consistent with each other.
\item \textbf{Uncalibrated scores and unvalidated thresholds.} Softmax outputs are not probabilities of disease; no calibration was fitted, and the application's 0.5 ``inconclusive'' threshold is not validated.
\item \textbf{Grad-CAM} is qualitative and does not validate predictions.
\item \textbf{No clinical validation.} The system has not been evaluated prospectively, by radiologists or in any clinical workflow, and must not be used for diagnosis.
\end{itemize}

\subsection{Future work}
The most important extensions are, in order: a causal test of the source shortcut (leave-one-source-out training for multi-source classes, and inference with regions outside the lungs masked); more seeds per configuration, so that runs can be compared as distributions, with paired tests for ROC-AUC differences beyond the descriptive McNemar test; a finer ablation of the preprocessing chain (single filters, filtering after resizing); at least one further backbone under the identical protocol; access to VinDr-CXR and JSRT so that Pleural Effusion and Nodule/Mass have more than one source; a second external cohort covering more classes, ideally a local hospital set collected with ethics approval; post-hoc calibration fitted on validation; and evaluation of the input gate on real non-radiograph images.
